\section{Two invariant fields giving the harmonic diagram}\label{har:sec:fields}

The preceding proof chooses a base from powers of one fixed line.
Here we give two alternatives. In the first, finite rank of the
linearized Picard group controls the horizontal divisors of the
determinant sections. In the second, equal-degree products force
their horizontal orders to be affine in the exponent. Both produce
a relative anticanonical section, rank-one invariant adjoint spaces,
and a vertical curvature kernel, so that
Propositions~\ref{har:prop:maximum} and~\ref{har:prop:euler} apply.

Throughout this section assume \(H^i(X,\mathcal O_X)=0\) for \(i>0\),
and let \(T\) act algebraically on the smooth connected projective
\(X\). Put \(n=\dim X\) and \(L=-K_X\), choose a compact-torus-invariant smooth
semipositive metric, and write \(\theta\) for its curvature and
\(r\) for its maximal rank. Apply Lemma~\ref{har:lem:forms}. If it
already produces an invariant anticanonical section, there is nothing
to construct; this includes the full-rank case. Otherwise write its
determinant sections as \(s_i\in H^0(X,N_i)^T\), where
\(N_i=M+a_iL\) and \(a_i\to\infty\).
Lemma~\ref{har:lem:null-divisors} gives their zero intersection with
\([\theta]^r[\omega]^{n-r-1}\). If \(r=0\), their effective zero
divisors have zero K\"ahler degree and hence vanish; a ratio of two
sections gives the conclusion directly. We therefore work with
\(0<r<n\).

\subsection{Invariant functions with curvature-null poles}\label{har:subsec:small}

Call a prime divisor small if \(\theta^r\) vanishes on its smooth
locus. An effective divisor has small support exactly when its
intersection with \([\theta]^r[\omega]^{n-r-1}\) vanishes.
Smallness persists under pullback to smooth birational
models and for exceptional divisors centered in small supports. To see
this even over a singular locus, resolve the original prime, take a
dominating component of the fiber product with the proposed source,
and test the pulled-back form on a smooth dominating space; continuity
extends its zero identity.

Let \(\mathcal K\) consist of zero and the invariant rational functions
with small pole support. It is a field: sums and products introduce no
new pole support, while zeros and poles of a rational function are
linearly equivalent effective divisors, so the zero support is small
whenever the pole support is. This also handles inversion. It is a
finitely generated intermediate field in \(\mathbb C(X)\). Explicitly,
take a transcendence basis of \(\mathcal K\), extend it to one of
\(\mathbb C(X)\), and note that adjoining the complementary independent
variables preserves the degrees of finite algebraic subextensions.
The finite degree of the full function field over the larger purely
transcendental field bounds those degrees, making the algebraic part
finite.

Use finitely many generators to map to a product of projective lines,
take the image \(Y_0\), and resolve equivariantly:
\(\pi_0:X_0\to X\), \(f_0:X_0\to Y_0\). The base action is trivial.
Indeterminacy lies in the small pole supports, so the resolution may
be chosen to be an isomorphism off them. The exceptional canonical
divisor \(E_0\) is small. A pullback of an ample line \(H_0\) on
\(Y_0\) has a small effective representative, built from the pole
divisors of the generators and exceptional corrections.

The choice of field now controls horizontal divisors. There are only
finitely many invariant small primes on \(X_0\) horizontal over
\(Y_0\). Indeed \(\operatorname{Pic}^0(X_0)=0\) by birational
invariance, and the kernel of the linearized Picard group over the
ordinary Picard group consists of characters. Thus the linearized
Picard group has finite rank. Infinitely many such primes would give,
after clearing torsion, a nonzero integral relation among their
natural linearized divisor lines. The corresponding quotient of
canonical divisor sections is an invariant rational function whose
divisor is that nonzero horizontal combination. Its poles downstairs
are small, so it belongs to \(\mathcal K=\mathbb C(Y_0)\).
But a function pulled back from \(Y_0\) has vertical divisor, a
contradiction.

The determinant divisors chosen above are small by
Lemma~\ref{har:lem:null-divisors}, so their horizontal supports belong
to this finite set. Order
the multiplicities along an infinite subsequence and take a ratio.
We obtain an invariant rational section \(s\) of \(b\pi_0^*L\),
\(b>0\), regular on the generic fiber. For every \(k\ge0\), both
\[
 (f_{0*}\mathcal O_{X_0}(kb\pi_0^*L))^T,
 \qquad (f_{0*}\mathcal O_{X_0}(E_0+kb\pi_0^*L))^T
\]
have rank one. They contain \(s^k\) and \(s_{E_0}s^k\). Any other
invariant generic section has its vertical poles cleared by a power
of \(H_0\). Its effective zero divisor is small, since its class has
zero intersection with the semipositive test forms, using smallness of
\(E_0,H_0\) and \(\theta^{r+1}=0\).
The ratio of two such sections therefore has small poles and is in
\(\mathcal K\). They are dependent over the base field, as required.

The kernel condition has a direct differential proof. Every level
divisor of a generating function is small. At a regular level point
in the rank-\(r\) open, its tangent hyperplane must contain
\(\ker\theta\): a hyperplane not containing that kernel would surject
onto the rank-\(r\) quotient, so \(\theta^r\) would not vanish on it.
Thus \(\ker\theta\subset\ker df_0\) on a nonempty ordinary open.
After an equidimensional normal model and a smooth resolution, all
these properties persist. Proposition~\ref{har:prop:maximum} gives the
continuous base metric and Proposition~\ref{har:prop:euler} gives an
invariant anticanonical section. No separate relative-algebraic-closure
assumption on the small-pole field is needed. The invariant adjoint
rank at \(k=0\) already forces geometric connectedness of the generic
fiber, as explained in Proposition~\ref{har:prop:maximum}.

\subsection{Ratios from a semigroup of determinant lines}\label{har:subsec:semigroup}

Retain the hypotheses and determinant sections fixed at the start of
the section. This time we impose all equal-degree ratios at once;
their relations will give a relative section without first listing
the horizontal primes. Form the semigroup of linearized line bundles generated by the
\(N_i\), and the field generated by ratios of invariant sections of
the same semigroup element. Choose a transcendence basis among these
ratios. Products put their numerators and denominators into one
invariant linear system of a semigroup element \(N_0\). Normalize its
image in the relative algebraic closure of its function field in
\(\mathbb C(X)\), and resolve the base to obtain \(Y\). Every original
ratio is algebraic over the chosen transcendence basis, so belongs
to this finite relative algebraic closure. The connected torus
acts trivially on it: it fixes the smaller field and has no nontrivial
finite quotient. Thus the base action is trivial.

After the Hilbert-family modification, take the normal equidimensional
main graph \(W\to Y\) and its smooth resolution \(V\). Write
\(\pi:V\to X\), \(f:V\to Y\), \(D_V=\pi^*L\), \(E=K_{V/X}\).
The pullback \(A\) of the image hyperplane bundle to \(Y\) is big,
and the common factor of the pulled-back coordinates is regular,
giving an invariant effective difference
\(\pi^*N_0-f^*A\). Hence for an ample \(H\) on \(Y\),
\begin{equation}\label{har:eq:absorb}
 a\pi^*N_0-f^*H\quad\hbox{is invariantly effective for some }a>0.
\end{equation}
The intersection identity of Lemma~\ref{har:lem:null-divisors} for
every \(N_i\), and hence \(N_0\), forces the intersection of \(f^*A\) with
\((\pi^*\theta)^r(\pi^*\omega)^{n-r-1}\) to vanish. Where the map to
the image has maximal rank this proves the kernel condition, by the
same positive-semidefinite linear algebra used above.

For every semigroup element \(N\) and integer \(e\ge0\),
\[
 \dim_{\mathbb C(Y)}H^0(V_\eta,\pi^*N+eE)^T\le1.
\]
If two generic sections were independent, twist their invariant direct
image by a sufficiently ample base line, obtain global invariant
sections, and absorb that twist with \eqref{har:eq:absorb}. Since
\(\pi_*\mathcal O_V(eE)=\mathcal O_X\), this gives invariant sections
of a semigroup element whose ratio is not in \(\mathbb C(Y)\),
contradicting its construction.

Let \(F_i\) be the horizontal divisor of \(\pi^*s_i\). For
\(a_i<a_j<a_h\), positive integers \(u,v\) with
\(ua_i+va_h=(u+v)a_j\) yield two nonzero invariant generic sections
of the same line. Their ratio is constant on the generic fiber, so
\(uF_i+vF_h=(u+v)F_j\). Consequently \(F_i\) is affine in \(a_i\).
Its slope is effective because \(F_i\ge0\) and \(a_i\to\infty\).
The ratio \(s=\pi^*s_2/\pi^*s_1\), of degree \(b=a_2-a_1>0\),
therefore has no horizontal poles. In fact
\[
 \dim_{\mathbb C(Y)}H^0(V_\eta,kbD_V+eE)^T=1
 \quad(k,e\ge0).
\]
The displayed section \(s^ks_E^e\) gives the lower bound. For the
upper bound choose \(i\) large enough that
\(F_i-k(F_2-F_1)\ge0\). Multiplication by \(\pi^*s_i/s^k\) embeds
the space into the preceding rank-at-most-one space for \(N_i+eE\).
If the base is a point, \(s\) is already globally regular. Otherwise
Propositions~\ref{har:prop:maximum} and \ref{har:prop:euler} finish this
route. The stronger all-\(e\) generic-rank statement is retained here;
only \(e=1\) is needed for the metric and transfer.
